\chapter{Conclusão}
\label{cap:conclusao}

Este trabalho investigou o uso de LLMs para simplificação e adaptação comunicacional de textos técnicos de saúde em um contexto mHealth. A motivação central partiu da dificuldade que conteúdos de saúde, frequentemente escritos em linguagem técnica, podem impor a usuários finais de aplicações digitais, especialmente quando esses conteúdos envolvem orientações de cuidado que precisam ser compreendidas de forma clara. Nesse cenário, o objetivo geral do trabalho foi investigar, implementar e avaliar uma solução baseada em LLMs para simplificação e adaptação de linguagem técnica em uma plataforma mHealth baseada em planos de cuidado, utilizando guardrails para apoiar a preservação do significado e das informações essenciais do material original. Esse objetivo foi desenvolvido a partir da fundamentação teórica, da análise de trabalhos relacionados, da proposta da solução, da implementação funcional da aplicação SimplyHealth e da avaliação com materiais de saúde previamente definidos, perfis simulados, indicador de legibilidade, análise qualitativa e verificação da atuação dos guardrails.

A proposta buscou enfrentar a dificuldade de compreensão de conteúdos técnicos de saúde por meio de um fluxo de simplificação textual orientado por perfis comunicacionais de pacientes. Nesse contexto, a Takere foi utilizada como ambiente de referência para situar a abordagem em um cenário mHealth baseado em planos de cuidado, sem que tenha sido realizada uma integração plena com a plataforma. A adaptação foi delimitada como comunicacional, e não clínica: os perfis simulados, definidos a partir de aspectos como idade, escolaridade e área de formação, foram utilizados apenas para ajustar vocabulário, tom, nível de explicação e forma de apresentação do conteúdo. Dessa forma, a LLM foi empregada como mecanismo de reescrita e adaptação linguística, não como fonte autônoma de conhecimento médico, nem como agente responsável por validar, recomendar ou modificar condutas clínicas.

Essa proposta foi materializada em uma implementação funcional, com frontend e backend, capaz de receber textos livres ou arquivos nos formatos PDF e TXT, processar o conteúdo original e gerar saídas voltadas à compreensão do usuário. A aplicação implementou a geração de texto simplificado, glossário, perguntas frequentes, chat contextual com entrada por voz, síntese de fala do conteúdo simplificado e apresentação do indicador de legibilidade Flesch-PT. Além disso, os prompts foram estruturados para orientar a tarefa da LLM e incorporar regras de preservação do conteúdo, enquanto os guardrails foram operacionalizados por meio de instruções no prompt, verificações determinísticas, avaliação por LLM e mecanismo de fallback. Essa implementação demonstra que o trabalho não se limitou a uma proposta conceitual, mas resultou em uma solução funcional voltada à investigação prática do uso de LLMs na simplificação de linguagem técnica em saúde.

Como trabalhos futuros, recomenda-se ampliar a avaliação para um conjunto maior de materiais de saúde, incluindo conteúdos com diferentes níveis de complexidade e criticidade. Também seria relevante realizar estudos com pacientes reais ou usuários finais, investigando compreensão textual, percepção de clareza e dificuldades de uso em situações mais próximas do contexto real. Outra possibilidade é envolver profissionais de saúde na revisão sistemática das saídas geradas, permitindo uma análise mais robusta da preservação das informações clínicas. Além disso, futuras versões da solução poderiam comparar diferentes modelos de linguagem, aprimorar os guardrails determinísticos e baseados em LLM, melhorar mecanismos de rastreabilidade das simplificações e avaliar uma integração mais próxima com a Takere. Assim, este trabalho contribui ao demonstrar, por meio de uma implementação funcional, que LLMs podem apoiar a adaptação comunicacional de textos técnicos de saúde, desde que utilizadas com escopo delimitado, mecanismos de controle e avaliação cuidadosa das informações preservadas.
